%% toc: P15 174 174.1,174.2,174.3,174.4
\chapter{Target Populations, Sampling, and Observational Comparisons}
\label{sec:p15-w02b-p15-075}

A comparison of treatment groups answers a question only after three things
are fixed: the units the question is about, the quantity to be learned about
them, and the probability law, if any, that connects the observed responses
to that quantity. \Cref{sec:p15-factorial-inference} works inside a fixed set
of recruited units. It treats estimable contrasts and design matrices
(\cref{sec:p15-factorial-estimability}), conditional, average and interaction
contrasts in the two-by-two design (\cref{sec:p15-factorial-two-by-two}),
replication, pure error and repeated measurements
(\cref{sec:p15-factorial-uncertainty}), completely randomized and blocked
designs (\cref{sec:p15-factorial-randomized-blocks}), full and fractional
factorials (\cref{sec:p15-factorial-full,sec:p15-factorial-fractions}), and
the unbiased assignment estimator with its exact variance
(\cref{thm:p15-factorial-assignment}). This chapter adds what that treatment
leaves open.

\begin{enumerate}
\item It separates the units recruited into a study from the target population
named by the question, defines an estimand for each, and states exactly which
premises connect them (\cref{sec:p15-w02b-p15-075-estimands}).
\item It turns the assignment variance identity into observable inference: a
conservative variance estimator, a proof that no unbiased one exists, a
finite-sample interval, a randomization test and the confidence set obtained by
inverting it, and the design-based analysis of blocked assignment
(\cref{sec:p15-w02b-p15-075-design}).
\item It distinguishes random sampling of units from a target population,
random assignment of treatments to units, and observational comparisons in
which the assignment law is neither controlled nor known
(\cref{sec:p15-w02b-p15-075-sampling,sec:p15-w02b-p15-075-observational}).
\end{enumerate}

The inputs are a question naming treatment levels, an outcome and a population;
a finite set of units, each with fixed potential responses; a sampling design
and an assignment law when the investigator controls them; and the realized
data. The conclusions are of three kinds: probability statements over a
known randomization; logical bounds when no law is known; and conditional
conclusions that hold under named premises which the data cannot check. One
hypothetical six-unit table (\cref{def:p15-w02b-p15-075-table}) is carried
through every section; it is an illustration, not a report of a conducted
study.

\section{Questions, target populations and estimands}
\label{sec:p15-w02b-p15-075-estimands}

Hern\'an and Robins describe a causal question as having three components: an
action, an outcome, and a specific population in which the action's effect on
the outcome is wanted \cite[ch.~2, p.~13]{HernanRobins2020WhatIf2026}. The
population is part of the question. A study of whether a checking instruction
improves accuracy on reasoning problems (\cref{def:p19-proposed-factorial-design})
may recruit forty volunteers from one course, while the question concerns the
students who will be taught with the instruction next year. The volunteers are
the recruited units; next year's students are the target population. The two
sets can overlap, be nested, or be disjoint.

\begin{hbdefinition}{def:p15-w02b-p15-075-populations}{Target population, recruited units and their estimands}
Fix a nonempty finite set \(\mathcal T\) of treatment levels and real weights
\((w_t)_{t\in\mathcal T}\). A \emph{target population} is a nonempty finite
set \(\mathcal U\) of \(M\) units. Each unit \(k\) carries real
\emph{potential responses} \(Y_k(t)\), \(t\in\mathcal T\), and a value
\(V_k\) of a covariate with finite value set \(\mathcal V\). The
\emph{recruited units} form a nonempty finite set \(\mathcal R\) of \(N\)
units carrying the same data; \(\mathcal R\) may or may not be a subset of
\(\mathcal U\). The \emph{unit contrast} of unit \(k\) is
\(\tau_k=\sum_{t}w_tY_k(t)\). The \emph{recruited-unit estimand} and the
\emph{target estimand} are
\[
 \tau=\frac1N\sum_{i\in\mathcal R}\tau_i,\qquad
 \tau_{\mathcal U}=\frac1M\sum_{k\in\mathcal U}\tau_k .
\]
For \(v\in\mathcal V\), let \(q_{\mathcal R}(v)\) and \(q_{\mathcal U}(v)\)
be the proportions of units with \(V=v\) in \(\mathcal R\) and in
\(\mathcal U\). When a proportion is positive, let \(\tau(v)\) and
\(\tau_{\mathcal U}(v)\) be the averages of the unit contrasts over the units
with \(V=v\) in \(\mathcal R\) and in \(\mathcal U\), respectively. The
estimand \(\tau\) is the one of \cref{def:p15-factorial-assignment}.
\index{target population}\index{estimand!recruited-unit}\index{estimand!target}
\end{hbdefinition}

Potential responses are fixed numbers. They are well defined only under the
premises recorded in \cref{def:p15-factorial-assignment}: stable treatment
versions, and no interference, so that one unit's response does not depend on
the treatments of other units. The estimands are finite averages of unit
contrasts. They are not parameters of an assumed probability model. In a
super-population formulation the same quantities become expectations under a
joint law of potential responses, as in
\cref{def:p16-binary-potential-outcomes}; Hern\'an and Robins call the
parameter of interest the estimand and the rule that maps data to a value the
estimator \cite[\S10.1, p.~138]{HernanRobins2020WhatIf2026}.

\begin{hbtheorem}{thm:p15-w02b-p15-075-estimand-gap}{proved=annex:p15-w02b-p15-075-estimand-gap}{Linking the recruited and target estimands}
Use the notation of \cref{def:p15-w02b-p15-075-populations}.
\begin{enumerate}
\item \emph{Composition and within-stratum terms.} Suppose every
\(v\) with \(q_{\mathcal U}(v)>0\) also has \(q_{\mathcal R}(v)>0\). Then
\[
 \tau_{\mathcal U}-\tau
 =\sum_{v:\,q_{\mathcal R}(v)>0}\bigl(q_{\mathcal U}(v)-q_{\mathcal R}(v)\bigr)\tau(v)
 +\sum_{v:\,q_{\mathcal U}(v)>0}q_{\mathcal U}(v)\bigl(\tau_{\mathcal U}(v)-\tau(v)\bigr).
\]
\item \emph{Stratum transport.} Under the same condition, if
\(\tau_{\mathcal U}(v)=\tau(v)\) for every \(v\) with \(q_{\mathcal U}(v)>0\),
then
\[
 \tau_{\mathcal U}=\sum_{v:\,q_{\mathcal U}(v)>0}q_{\mathcal U}(v)\tau(v).
\]
If in addition \(q_{\mathcal U}=q_{\mathcal R}\), then
\(\tau_{\mathcal U}=\tau\).
\item \emph{Nested recruitment.} Suppose \(\mathcal R\subseteq\mathcal U\)
and \(N<M\). Then
\(\tau_{\mathcal U}=(N/M)\tau+(1-N/M)\tau_{\mathcal U\setminus\mathcal R}\),
where \(\tau_{\mathcal U\setminus\mathcal R}\) is the average unit contrast of
the unrecruited units. If the only information about the unrecruited units is
that each unit contrast lies in a known interval \([c_-,c_+]\), then the set
of values of \(\tau_{\mathcal U}\) compatible with the recruited units'
potential responses is exactly
\[
 \Bigl[\tfrac NM\tau+\bigl(1-\tfrac NM\bigr)c_-,\ \
       \tfrac NM\tau+\bigl(1-\tfrac NM\bigr)c_+\Bigr].
\]
\item \emph{Disjoint recruitment.} If \(\mathcal R\cap\mathcal U=\varnothing\)
and the only information about \(\mathcal U\) is that each unit contrast lies in
\([c_-,c_+]\), then every value in \([c_-,c_+]\) is compatible with the
recruited units' potential responses, whatever they are.
\end{enumerate}
\end{hbtheorem}

Parts~(3) and~(4) concern the potential responses themselves, not estimates:
even complete knowledge of every recruited unit's potential responses leaves the
target estimand undetermined unless a premise links the two sets. Three kinds
of premise are used in practice. The first is identity, \(\mathcal R=\mathcal U\).
The second is a probability sample from \(\mathcal U\) with a known design
(\cref{sec:p15-w02b-p15-075-sampling}). The third is the stratum premise of
part~(2), together with known target proportions \(q_{\mathcal U}\). The
stratum premise is a claim about units that were never observed. Hern\'an and
Robins call extrapolation of this kind transportability. They note that it
requires the distribution of effect modifiers, the definition of treatment and
the pattern of interference to be suitably similar in the two populations, and
that a reweighted stratum average carries no guarantee when unmeasured
effect modifiers are distributed differently \cite[\S4.3, p.~52; Fine
Point~4.2, p.~54]{HernanRobins2020WhatIf2026}. Part~(1) makes the two ways the
premise can fail explicit: a composition term, which known proportions
correct, and a within-stratum term, which no recruited data can measure.

\begin{hbmethod}{def:p15-w02b-p15-075-estimand-procedure}{Fixing the estimand before the design}
\begin{enumerate}
\item State the treatment levels \(\mathcal T\), the outcome, and the contrast
weights \(w_t\). The weights of \cref{def:p15-factorial-coding} give the
conditional, average and interaction contrasts of a two-by-two design.
\item Name the target population \(\mathcal U\) and the recruited units
\(\mathcal R\), and state whether \(\mathcal R=\mathcal U\),
\(\mathcal R\subseteq\mathcal U\), or neither.
\item Choose the estimand: \(\tau\) if the conclusion concerns only the
recruited units, \(\tau_{\mathcal U}\) otherwise.
\item For \(\tau_{\mathcal U}\), name the linking premise: identity, a known
probability sampling design, or the stratum premise of
\cref{thm:p15-w02b-p15-075-estimand-gap}(2) with stated covariates \(V\) and
known proportions \(q_{\mathcal U}\). Report the target-level conclusion as
conditional on that premise.
\item Only then choose the assignment law and the analysis
(\cref{sec:p15-w02b-p15-075-design}).
\end{enumerate}
\end{hbmethod}

\section{Design-based inference for the recruited-unit contrast}
\label{sec:p15-w02b-p15-075-design}

In design-based inference the potential responses are fixed and the only
randomness is the assignment, whose law the investigator chose. Hern\'an and
Robins distinguish this target, inference restricted to the units actually
randomized, from inference about a super-population from which those units are
imagined to have been sampled \cite[\S10.2, pp.~141--142]{HernanRobins2020WhatIf2026}.
This section treats the first; \cref{sec:p15-w02b-p15-075-sampling} treats the
second when a real sampling design exists.

\begin{hbdefinition}{def:p15-w02b-p15-075-complete}{Complete randomization and its population quantities}
Let \(\mathcal R=\{1,\ldots,N\}\), let \(|\mathcal T|\geq2\), and fix cell
sizes \(n_t\geq1\) with \(\sum_tn_t=N\). \emph{Complete randomization}
draws the assignment \(Z=(Z_1,\ldots,Z_N)\) uniformly from the
\(N!/\prod_tn_t!\) assignments that place exactly \(n_t\) units in cell
\(t\); then \(\pi_i(t)=n_t/N\) in \cref{def:p15-factorial-assignment}. The
observed response is \(Y_i=Y_i(Z_i)\). Write
\[
\begin{gathered}
 \bar y(t)=\frac1N\sum_iY_i(t),\qquad
 S_t^2=\frac1{N-1}\sum_i\bigl(Y_i(t)-\bar y(t)\bigr)^2,\\
 S_w^2=\frac1{N-1}\sum_i(\tau_i-\tau)^2 .
\end{gathered}
\]
Here \(S_w^2\) is the variance of the unit contrasts across the recruited
units. Within cell \(t\), let \(\overline Y_t\) be the mean observed
response and, when \(n_t\geq2\),
\(s_t^2=(n_t-1)^{-1}\sum_{i:Z_i=t}(Y_i-\overline Y_t)^2\). The estimator is
\(\widehat\tau=\sum_tw_t\overline Y_t\), the form taken by the estimator of
\cref{thm:p15-factorial-assignment} under complete randomization.
\end{hbdefinition}

\begin{hbtheorem}{thm:p15-w02b-p15-075-neyman}{proved=annex:p15-w02b-p15-075-neyman}{Exact variance and a conservative variance estimator}
Under complete randomization:
\begin{enumerate}
\item \(\mathbb E[\widehat\tau]=\tau\) and
\[
 \operatorname{Var}(\widehat\tau)
 =\sum_t\frac{w_t^2S_t^2}{n_t}-\frac{S_w^2}{N}.
\]
\item If \(n_t\geq2\) whenever \(w_t\neq0\), then
\(\mathbb E[s_t^2]=S_t^2\) for each such \(t\), and
\(\widehat V=\sum_{t:\,w_t\neq0}w_t^2s_t^2/n_t\) satisfies
\[
 \mathbb E[\widehat V]=\operatorname{Var}(\widehat\tau)+\frac{S_w^2}{N}
 \ \geq\ \operatorname{Var}(\widehat\tau),
\]
with equality if and only if every unit has the same contrast,
\(\tau_i=\tau\).
\item Suppose two distinct levels \(t\neq u\) have \(w_tw_u\neq0\). Then no
function of the observed data \((Z,(Y_i)_{i\leq N})\) has expectation
\(\operatorname{Var}(\widehat\tau)\) for every table of potential responses.
\end{enumerate}
\end{hbtheorem}

The excess \(S_w^2/N\) in part~(2) is the price of never observing two
potential responses of the same unit. Part~(3) shows that the price cannot be
removed by a cleverer unbiased estimator: the variance contains products
\(Y_i(t)Y_i(u)\) that no single assignment reveals. Hern\'an and Robins call an
interval \emph{conservative} when it covers the estimand more often than its
nominal level \cite[\S10.1, p.~139]{HernanRobins2020WhatIf2026}; an interval
built from \(\widehat V\) errs in that direction to the extent that
\(\widehat V\) is near its expectation.

\begin{hbtheorem}{thm:p15-w02b-p15-075-chebyshev}{proved=annex:p15-w02b-p15-075-chebyshev}{A finite-sample interval for the recruited-unit contrast}
Under complete randomization, let \(0<\eta<1\).
\begin{enumerate}
\item If a number \(B\) known before the data satisfies
\(B\geq\operatorname{Var}(\widehat\tau)\), then
\[
 \Pr\Bigl(\,|\widehat\tau-\tau|\leq\sqrt{B/\eta}\,\Bigr)\geq1-\eta .
\]
\item If every potential response lies in a known interval \([m_-,m_+]\)
and \(h=m_+-m_-\), then
\[
 B=\frac{N}{N-1}\cdot\frac{h^2}{4}\sum_t\frac{w_t^2}{n_t}
\]
satisfies \(B\geq\operatorname{Var}(\widehat\tau)\).
\end{enumerate}
\end{hbtheorem}

The interval of part~(1) is valid at every sample size, but it is wide. Its
coverage argument needs a bound fixed in advance. Replacing \(B\) by the random
\(\widehat V\) gives the interval most often reported. Its coverage is a
large-sample property that requires a central limit theorem for
complete randomization, and this chapter does not prove one. The randomization
test below gives exact statements of a different kind.

\begin{hbdefinition}{def:p15-w02b-p15-075-sharp}{Sharp hypotheses and randomization p-values}
A \emph{sharp hypothesis} \(H\) is a rule that, from the realized assignment
\(z\) and observed responses \(y\), determines every potential response
\(Y_i(t)\). Two examples are used. The hypothesis of no effect for any unit,
\(H_0\), states \(Y_i(t)=Y_i(u)\) for all \(i,t,u\); it gives
\(Y_i(t)=y_i\). For two levels \(\mathcal T=\{0,1\}\) and real \(\delta\), the
constant-effect hypothesis \(H_\delta\) states \(Y_i(1)=Y_i(0)+\delta\) for
all \(i\); it gives \(Y_i(0)=y_i-\delta z_i\) and
\(Y_i(1)=Y_i(0)+\delta\). For each assignment \(z'\), let
\(y^H(z')=(Y_i(z'_i))_i\) be the responses that \(H\) says would have been
observed under \(z'\). Fix a real statistic \(T_H(z',y')\) before seeing the
data. With \(Z'\) drawn from the design law, the \emph{randomization
p-value} is
\[
 p_H=\Pr\bigl(T_H(Z',y^H(Z'))\geq T_H(z,y)\bigr),
\]
computed by enumerating the assignments or by an exact count.
\end{hbdefinition}

\begin{hbtheorem}{thm:p15-w02b-p15-075-randomization-test}{proved=annex:p15-w02b-p15-075-randomization-test}{Validity of the randomization test and its inversion}
Let the assignment law be known and have finite support.
\begin{enumerate}
\item If the sharp hypothesis \(H\) is true, then
\(\Pr(p_H\leq\alpha)\leq\alpha\) for every \(\alpha\in[0,1]\).
\item For two levels, suppose the constant-effect model
\(Y_i(1)=Y_i(0)+\delta_0\) holds for some \(\delta_0\). Let
\(C=\{\delta\in\mathbb R:p_{H_\delta}>\alpha\}\). Then
\(\Pr(\delta_0\in C)\geq1-\alpha\). Under this model
\(\tau=\delta_0\) for the weights \(w_0=-1\), \(w_1=1\).
\end{enumerate}
\end{hbtheorem}

The hypothesis \(H_0\) is stronger than \(\tau=0\). Part~(1) controls the
rejection rate only when every unit's contrast is zero. If unit contrasts
differ and average to zero, the theorem gives no bound. Likewise the set \(C\)
has guaranteed coverage only under the constant-effect model. With
heterogeneous effects it is a set of constant shifts not rejected by the data;
it is not a confidence set for \(\tau\).

\begin{hbdefinition}{def:p15-w02b-p15-075-blocked}{Blocked complete randomization}
Partition \(\mathcal R\) into blocks \(b\in\mathcal B\) of sizes \(N_b\).
In block \(b\) fix cell sizes \(n_{tb}\geq1\) with
\(\sum_tn_{tb}=N_b\), and completely randomize within each block,
independently across blocks. Write \(\tau_b\), \(S_{tb}^2\),
\(S_{wb}^2\), \(\overline Y_{tb}\) and \(s_{tb}^2\) for the quantities of
\cref{def:p15-w02b-p15-075-complete} computed inside block \(b\). The
blocked estimator is
\[
 \widehat\tau_{\mathcal B}=\sum_b\frac{N_b}{N}\sum_tw_t\overline Y_{tb}.
\]
\end{hbdefinition}

\begin{hbtheorem}{thm:p15-w02b-p15-075-blocked}{proved=annex:p15-w02b-p15-075-blocked}{What blocked assignment adds}
Under blocked complete randomization:
\begin{enumerate}
\item \(\tau=\sum_b(N_b/N)\tau_b\) and
\(\mathbb E[\widehat\tau_{\mathcal B}]=\tau\).
\item \(\operatorname{Var}(\widehat\tau_{\mathcal B})
=\sum_b(N_b/N)^2V_b\), where
\(V_b=\sum_tw_t^2S_{tb}^2/n_{tb}-S_{wb}^2/N_b\).
\item If \(n_{tb}\geq2\) whenever \(w_t\neq0\), then
\(\widehat V_{\mathcal B}=\sum_b(N_b/N)^2\sum_{t:\,w_t\neq0}w_t^2s_{tb}^2/n_{tb}\)
has expectation
\(\operatorname{Var}(\widehat\tau_{\mathcal B})
+\sum_b(N_b/N)^2S_{wb}^2/N_b\).
\end{enumerate}
\end{hbtheorem}

The blocked variance depends only on spread inside blocks. Variation of the
potential responses between blocks does not enter it. That is the gain from
blocking on a covariate that predicts the response. Part~(3) needs two
assigned units per cell in each block. With one unit per treatment in a block,
as in the additive analysis of \cref{thm:p15-factorial-blocks}, the design
alone supplies no variance estimate, and the model-based residual uses the
additional no-interaction assumption stated there. Repeated measurements on one
unit remain one assigned unit (\cref{thm:p15-factorial-replication}). They
change the responses \(Y_i(t)\), for example into averages, but not \(N\) or
\(n_t\).

\section{Random sampling from a target population}
\label{sec:p15-w02b-p15-075-sampling}

Random assignment decides which potential response of a recruited unit is
observed. Random sampling decides which units of the target population are
observed at all. The two are separate randomizations with separate
consequences.

\begin{hbdefinition}{def:p15-w02b-p15-075-sampling-design}{Sampling designs and inclusion probabilities}
A \emph{sampling design} on the target population \(\mathcal U\) is a
probability law \(p\) on the subsets of \(\mathcal U\), fixed by the
investigator and not depending on any unit's responses. The sample
\(\mathcal S\) has \(\Pr(\mathcal S=s)=p(s)\). The \emph{inclusion
probabilities} are \(\rho_k=\Pr(k\in\mathcal S)\) and
\(\rho_{kl}=\Pr(k\in\mathcal S,\,l\in\mathcal S)\), with
\(\rho_{kk}=\rho_k\). \emph{Simple random sampling} of size \(n\) is the
uniform law on the subsets of size \(n\). \emph{Stratified simple random
sampling} applies independent simple random samples inside the strata of a
covariate. For real values \(y_k\), write
\(\bar y_{\mathcal U}=M^{-1}\sum_ky_k\) and, when \(M\geq2\),
\(S_{\mathcal U}^2=(M-1)^{-1}\sum_k(y_k-\bar y_{\mathcal U})^2\).
\end{hbdefinition}

\begin{hbtheorem}{thm:p15-w02b-p15-075-inclusion}{proved=annex:p15-w02b-p15-075-inclusion}{Inverse-inclusion estimation}
Let \(y_k\) be fixed real values on \(\mathcal U\).
\begin{enumerate}
\item If \(\rho_k>0\) for every \(k\), then
\(\widehat y=M^{-1}\sum_{k\in\mathcal S}y_k/\rho_k\) satisfies
\(\mathbb E[\widehat y]=\bar y_{\mathcal U}\), and
\[
 \operatorname{Var}(\widehat y)
 =\frac1{M^2}\sum_{k,l\in\mathcal U}
  (\rho_{kl}-\rho_k\rho_l)\frac{y_ky_l}{\rho_k\rho_l}.
\]
\item If in addition \(\rho_{kl}>0\) for all \(k,l\), then
\[
 \widehat v=\frac1{M^2}\sum_{k,l\in\mathcal S}
  \frac{\rho_{kl}-\rho_k\rho_l}{\rho_{kl}}\cdot\frac{y_ky_l}{\rho_k\rho_l}
\]
satisfies \(\mathbb E[\widehat v]=\operatorname{Var}(\widehat y)\).
\item For simple random sampling of size \(n\), \(\rho_k=n/M\),
\(\widehat y\) is the sample mean, and for \(M\geq2\)
\(\operatorname{Var}(\widehat y)=(1/n-1/M)S_{\mathcal U}^2\).
\item If \(\rho_k=0\) for some \(k\), then no function of
\((\mathcal S,(y_l)_{l\in\mathcal S})\) has expectation
\(\bar y_{\mathcal U}\) for every choice of the values \(y\).
\end{enumerate}
\end{hbtheorem}

Part~(3) agrees with the finite-population variance of
\cref{thm:p15-without-replacement-convex}. Part~(4) is the sampling analogue of
the nested-recruitment bound: a unit the design can never select is a unit
about which the sample is silent.

\begin{hbtheorem}{thm:p15-w02b-p15-075-two-stage}{proved=annex:p15-w02b-p15-075-two-stage}{Sampling followed by assignment}
Draw \(\mathcal S\) by a known sampling design with every \(\rho_k>0\).
Given \(\mathcal S=s\), assign the units of \(s\) to levels by a known law with
\(\pi_k(t\mid s)=\Pr(Z_k=t\mid\mathcal S=s)>0\) for \(k\in s\). Suppose the
joint law of \((\mathcal S,Z)\) does not depend on the potential responses, and
observe \(Y_k=Y_k(Z_k)\) for \(k\in\mathcal S\).
\begin{enumerate}
\item The statistic
\[
 \widehat\tau_{\mathcal U}=\frac1M\sum_{k\in\mathcal S}\sum_t
 \frac{w_t\mathbf1_{\{Z_k=t\}}Y_k}{\rho_k\,\pi_k(t\mid\mathcal S)}
\]
satisfies \(\mathbb E[\widehat\tau_{\mathcal U}]=\tau_{\mathcal U}\).
\item Suppose \(\mathcal S\) is a simple random sample of size \(n<M\), and
given \(\mathcal S\) the assignment is a complete randomization with cell sizes
\(n_t\geq1\), \(\sum_tn_t=n\). Then
\(\widehat\tau_{\mathcal U}=\sum_tw_t\overline Y_t\) and
\[
 \operatorname{Var}(\widehat\tau_{\mathcal U})
 =\sum_t\frac{w_t^2S_{\mathcal U,t}^2}{n_t}-\frac{S_{\mathcal U,w}^2}{M},
\]
where \(S_{\mathcal U,t}^2\) and \(S_{\mathcal U,w}^2\) are the quantities of
\cref{def:p15-w02b-p15-075-complete} computed over all \(M\) units of
\(\mathcal U\). If \(n_t\geq2\) whenever \(w_t\neq0\), the statistic
\(\widehat V\) of \cref{thm:p15-w02b-p15-075-neyman} has expectation
\(\operatorname{Var}(\widehat\tau_{\mathcal U})+S_{\mathcal U,w}^2/M\).
\end{enumerate}
\end{hbtheorem}

The same arithmetic \(\sum_tw_t\overline Y_t\) and the same \(\widehat V\)
thus serve two estimands with two variances. For the recruited units the
conservative excess is \(S_w^2/N\). For a target population sampled at random
it is \(S_{\mathcal U,w}^2/M\), which is small when \(M\) is large compared with
the sample. This is the finite form of the observation by Hern\'an and Robins
that randomizing and then sampling, or sampling and then randomizing, leads to
the same super-population analysis \cite[\S10.2, pp.~141--142]{HernanRobins2020WhatIf2026}.
They also point out that the super-population is usually a fiction. Most
studies do not sample at random from any large population, so an interval
computed as if they did is a statement about what would hold had they done so
\cite[\S10.3, pp.~142--144]{HernanRobins2020WhatIf2026}. Without a real
sampling design, the link from \(\tau\) to \(\tau_{\mathcal U}\) is one of the
premises of \cref{thm:p15-w02b-p15-075-estimand-gap}.

\begin{hbcorollary}{thm:p15-w02b-p15-075-licenses}{proved=annex:p15-w02b-p15-075-licenses}{What each randomization licenses}
\begin{enumerate}
\item \emph{Assignment without sampling.} Complete randomization of
\(\mathcal R\) gives an unbiased estimator of \(\tau\)
(\cref{thm:p15-w02b-p15-075-neyman}). If \(\mathcal R\subsetneq\mathcal U\)
and nothing links the unrecruited units, every value in the interval of
\cref{thm:p15-w02b-p15-075-estimand-gap}(3) remains compatible with the
recruited units' potential responses.
\item \emph{Sampling without assignment.} Suppose each target unit has a
two-level treatment \(Z_k\in\{0,1\}\) fixed by a process outside the
investigator's control, both levels occur in \(\mathcal U\), and the sample is
drawn by a design with every \(\rho_k>0\). The inverse-inclusion estimator of
the mean observed response \(M^{-1}\sum_kY_k(Z_k)\) is unbiased. However, the
population difference of observed means between \(Z=1\) and \(Z=0\) equals
\(\tau_1+\sigma\), the decomposition of
\cref{thm:p15-w02b-p15-075-observational}(1) computed over \(\mathcal U\).
Neither \(\tau_1-\tau_{\mathcal U}\) nor \(\sigma\) need vanish. Random sampling reproduces the population's
association; it does not convert association into the causal contrast.
\end{enumerate}
\end{hbcorollary}

\section{Observational comparisons}
\label{sec:p15-w02b-p15-075-observational}

An observational study is one in which the investigator observes and records
the data but does not carry out the action of interest
\cite[ch.~3, p.~27]{HernanRobins2020WhatIf2026}. In design terms, the
assignment vector is produced by a mechanism the investigator neither controls
nor knows. Its law may depend on the units' potential responses, for example
when the units who would fare worst without treatment are the ones selected for
it.

\begin{hbdefinition}{def:p15-w02b-p15-075-observational}{An observational comparison}
Let \(\mathcal T=\{0,1\}\), \(w_0=-1\), \(w_1=1\), so that
\(\tau=N^{-1}\sum_i(Y_i(1)-Y_i(0))\). An \emph{observational comparison} records
an assignment \(z\in\{0,1\}^N\) with \(n_1=\sum_iz_i\) and \(n_0=N-n_1\) both
positive, and the responses \(y_i=Y_i(z_i)\). The law that produced \(z\) is
not known; it may depend on the potential responses and on features of the
units that were not recorded. The \emph{naive contrast} is
\(D(z)=\overline Y_1-\overline Y_0\).
\end{hbdefinition}

\begin{hbtheorem}{thm:p15-w02b-p15-075-observational}{proved=annex:p15-w02b-p15-075-observational}{The naive contrast without a known assignment law}
In an observational comparison:
\begin{enumerate}
\item \(D(z)=\tau_1(z)+\sigma(z)\), where
\[
\begin{aligned}
 \tau_1(z)&=\frac1{n_1}\sum_{i:z_i=1}\bigl(Y_i(1)-Y_i(0)\bigr),\\
 \sigma(z)&=\frac1{n_1}\sum_{i:z_i=1}Y_i(0)-\frac1{n_0}\sum_{i:z_i=0}Y_i(0).
\end{aligned}
\]
The first term is the average contrast among the treated units. The second,
the \emph{selection term}, compares the treated and untreated units' responses
under the control level.
\item Suppose every unit has the same contrast \(c\), so \(\tau=c\), and the
values \(Y_i(0)\) are not all equal. If the mechanism selects for treatment the
\(n_1\) units with the smallest values of \(Y_i(0)\), ties broken in any fixed
way, then \(D(z)<\tau\) with probability one. If it selects the \(n_1\) units
with the largest values, then \(D(z)>\tau\).
\item If every potential response lies in a known interval \([m_-,m_+]\),
\(h=m_+-m_-\), then the values of \(\tau\) compatible with the recorded
\((z,y)\) form exactly the interval \([L(z),U(z)]\), where
\[
\begin{aligned}
 L(z)&=\frac1N\Bigl(\sum_{i:z_i=1}y_i-\sum_{i:z_i=0}y_i+n_0m_--n_1m_+\Bigr),\\
 U(z)&=\frac1N\Bigl(\sum_{i:z_i=1}y_i-\sum_{i:z_i=0}y_i+n_0m_+-n_1m_-\Bigr),
\end{aligned}
\]
and \(U(z)-L(z)=h\). The same interval is the logical range for a randomized
realization with the same \((z,y)\).
\end{enumerate}
\end{hbtheorem}

Part~(3) locates the difference between a randomized and an observational
comparison. It does not lie in the realized table, which constrains \(\tau\)
equally in both cases. It lies in the law: complete randomization makes
\(D\) unbiased with the variance of \cref{thm:p15-w02b-p15-075-neyman}, while an
unknown law supplies no probability statement about \(D-\tau\). Hern\'an and
Robins list the three identifiability conditions under which an observational
study can be analysed as a conditionally randomized experiment: consistency,
exchangeability conditional on measured covariates, and positivity
\cite[\S3.1, p.~28]{HernanRobins2020WhatIf2026}. The next result is their finite
design-based counterpart.

\begin{hbtheorem}{thm:p15-w02b-p15-075-stratified}{proved=annex:p15-w02b-p15-075-stratified}{Stratification under an assignment premise}
In an observational comparison, let a recorded covariate partition
\(\mathcal R\) into strata \(b\in\mathcal B\) of sizes \(N_b\), with realized
treated counts \(n_{1b}\) and untreated counts \(n_{0b}=N_b-n_{1b}\).
Consider two premises.
\begin{itemize}
\item[(P)] Conditional on the stratum counts \((n_{1b})_b\), every assignment
with those counts had the same probability.
\item[(Q)] \(1\leq n_{1b}\leq N_b-1\) for every \(b\).
\end{itemize}
\begin{enumerate}
\item Under (P) and (Q), conditional on the counts,
\(\widehat\tau_{\rm str}=\sum_b(N_b/N)(\overline Y_{1b}-\overline Y_{0b})\)
has expectation \(\tau\) and the variance of
\cref{thm:p15-w02b-p15-075-blocked}(2). If \(n_{1b},n_{0b}\geq2\) for every
\(b\), the estimator of \cref{thm:p15-w02b-p15-075-blocked}(3) is
conservative for that variance.
\item Under (P) alone, \(D(z)\) need not have expectation \(\tau\), even when
the conditional law is that of (P).
\item No recorded \((z,x,y)\) excludes (P), and none excludes its failure: the
law of (P) and the law placing all its mass on the realized \(z\) both give
the realized assignment positive probability, and the set of potential-response
tables compatible with \((z,y)\) does not depend on the law.
\end{enumerate}
\end{hbtheorem}

Premise (P) is the finite form of conditional exchangeability, and (Q) is
positivity. In the probabilistic formulation of
\cref{def:p16-binary-potential-outcomes}, the same two premises with
consistency identify the average treatment effect by the standardized formula
of \cref{thm:p16-po-ate-identification}. That formula has the form of
\(\widehat\tau_{\rm str}\), with population stratum weights in place of
\(N_b/N\). When positivity fails, so that some stratum has only treated or only
untreated units, \(\widehat\tau_{\rm str}\) is undefined; restricting the
estimand to the strata where (Q) holds changes the estimand. Part~(3) is why
an observational conclusion is conditional on (P): the premise concerns the
unrecorded mechanism, not the data.

\begin{hbwarning}{def:p15-w02b-p15-075-confounded-warning}{A confounded comparison can reverse the sign}
If the units chosen for treatment would have had lower responses without it,
the selection term of \cref{thm:p15-w02b-p15-075-observational}(1) is negative.
It can outweigh a positive effect. In the table below, every unit's contrast is
nonnegative and \(\tau=3/2\), yet selecting the three units with the lowest
control responses gives \(D=-5/3\). Adjusting for a covariate removes this term
only under premise (P) for that covariate.
\end{hbwarning}

\section{One table under three designs}
\label{sec:p15-w02b-p15-075-example}

\begin{hbexample}{def:p15-w02b-p15-075-table}{A hypothetical six-unit table}
Six recruited units have a binary covariate \(X\), with \(X=1\) marking a
severe condition, and potential responses on a scale from \(0\) to \(10\),
higher being better:
\[
\begin{array}{c|cccccc}
 i&1&2&3&4&5&6\\\hline
 X_i&1&1&1&0&0&0\\
 Y_i(0)&1&2&3&4&5&6\\
 Y_i(1)&3&3&4&6&5&9\\
 \tau_i&2&1&1&2&0&3
\end{array}
\]
The contrast weights are \(w_0=-1\), \(w_1=1\). The numbers are constructed
for illustration; no study produced them.
\end{hbexample}

\begin{hbtheorem}{thm:p15-w02b-p15-075-worked}{proved=annex:p15-w02b-p15-075-worked}{The worked calculations}
For the table of \cref{def:p15-w02b-p15-075-table}:
\begin{enumerate}
\item \emph{Estimands.} \(\tau=3/2\), \(S_1^2=26/5\), \(S_0^2=7/2\),
\(S_w^2=11/10\). Within strata, \(\tau(1)=4/3\) and \(\tau(0)=5/3\).
\item \emph{Complete randomization}, \(n_1=n_0=3\).
\(\operatorname{Var}(\widehat\tau)=163/60\) and
\(\mathbb E[\widehat V]=29/10=163/60+11/60\).
For the realized treated set \(\{1,4,5\}\), the observed responses are
\((3,2,3,6,5,6)\), \(\widehat\tau=1\) and \(\widehat V=20/9\).
\item \emph{Randomization test of \(H_0\).} With the statistic
\(\overline Y_1-\overline Y_0\), the twenty null values and their
multiplicities are
\[
\begin{array}{c|cccccccc}
 \text{value}&-3&-5/3&-1&-1/3&1/3&1&5/3&3\\\hline
 \text{multiplicity}&1&2&5&2&2&5&2&1
\end{array}.
\] The one-sided p-value is \(2/5\); with the statistic
\(|\overline Y_1-\overline Y_0|\) the two-sided p-value is \(4/5\).
\item \emph{Inversion.} With the statistic
\(|\overline Y_1-\overline Y_0-\delta|\) for \(H_\delta\) and
\(\alpha=1/10\), the set \(C\) is the closed interval \([-3,4]\).
\item \emph{Chebyshev interval.} With \([m_-,m_+]=[0,10]\) and
\(\eta=1/5\), \(B=20\) and the interval is \([-9,11]\).
\item \emph{Observational selection.} If the three units with the lowest
\(Y_i(0)\), namely \(\{1,2,3\}\), are treated, then \(D=-5/3\),
\(\tau_1=4/3\), \(\sigma=-3\), and the logical range of \(\tau\) given the
observed data is \([-35/6,25/6]\).
\item \emph{Stratified mechanism.} If two of the three severe units and one of
the three others are treated, each choice uniform and independent across strata,
then \(\mathbb E[D]=4/9\), while \(\widehat\tau_{\rm str}\) has expectation
\(3/2\) and variance \(23/18\).
\item \emph{Target population.} If \(q_{\mathcal U}(1)=1/5\) and the stratum
premise holds, then \(\tau_{\mathcal U}=8/5\). The composition term is
\(1/10\), and the within-stratum term is zero. If instead the six units are part
of a target population of \(M=30\) units whose unit contrasts are known only to
lie in \([-10,10]\), then \(\tau_{\mathcal U}\) can be any value in
\([-77/10,83/10]\).
\end{enumerate}
\end{hbtheorem}

The worked numbers separate the conclusions. The randomized estimate
\(\widehat\tau=1\) differs from \(\tau=3/2\) by chance alone, and the test,
the inverted set and the Chebyshev interval each quantify that chance under
their own premises. The inverted set \([-3,4]\) contains \(\tau\), but the
table violates the constant-effect model, so this coverage is not one the
theorem guarantees. The Chebyshev interval is longer than the logical range
\([-10,10]\) of a contrast on this scale. At six units the guaranteed
finite-sample interval carries no information; its half-width shrinks like
\(N^{-1/2}\) at fixed cell proportions. The observational selection reverses
the sign. The stratified mechanism leaves the naive contrast biased by
\(4/9-3/2=-19/18\), and the stratified estimator removes the bias only because
premise (P) holds by construction here. Finally, even exact knowledge of all
twelve potential responses leaves \(\tau_{\mathcal U}\) undetermined over an
interval of length sixteen, unless a sampling design or a transport premise is
added.

\section{Limitations and relationships}
\label{sec:p15-w02b-p15-075-relationships}

\paragraph{What the results require.}
Every result here assumes well-defined potential responses: stable treatment
versions, no interference, and consistency between the assigned level and the
recorded response. Differential loss of responses, noncompliance and
measurement error break the observation model, as noted after
\cref{thm:p15-factorial-assignment}. The variance results of
\cref{thm:p15-w02b-p15-075-neyman,thm:p15-w02b-p15-075-blocked} are exact
for the stated designs. The Chebyshev interval is exact at every size but wide.
The usual normal-approximation interval is not established here. Randomization
tests are exact only for sharp hypotheses.

\paragraph{What the results do not license.}
An unbiased assignment estimate concerns \(\tau\), an average over recruited
units; it is not a claim about any single unit, since \(S_w^2\) and the
individual contrasts are not identified. A target-population conclusion needs
one of the linking premises of \cref{thm:p15-w02b-p15-075-estimand-gap}. A
mechanism claim needs further premises again, as
\cref{thm:p15-factorial-mechanism} shows for the factorial surface.

\paragraph{Relationships.}
\Cref{thm:p15-w02b-p15-075-neyman} specializes
\cref{thm:p15-factorial-assignment} to complete randomization and makes its
variance observable up to the excess \(S_w^2/N\). Blocked assignment
(\cref{thm:p15-w02b-p15-075-blocked}) is the design-based counterpart of the
additive block analysis of \cref{thm:p15-factorial-blocks}. The general
weights \(w_t\) include every contrast of \cref{def:p15-factorial-coding}, so
the conditional, average and interaction contrasts of a two-by-two design
inherit the variance, the conservative estimator and the interval. In a
two-by-two design the unit contrast of the difference of differences is
\(Y_i(11)-Y_i(01)-Y_i(10)+Y_i(00)\), and its variance across units enters
\(S_w^2\). The super-population identification results of
\cref{thm:p16-po-ate-identification,thm:p16-po-att-identification} correspond
to \cref{thm:p15-w02b-p15-075-stratified}. The finite-population variance of
\cref{thm:p15-without-replacement-convex} is the special case behind
\cref{thm:p15-w02b-p15-075-inclusion}(3) and
\cref{thm:p15-w02b-p15-075-two-stage}(2).

\paragraph{Sources.}
The distinctions between question, population, estimand and estimator, between
randomization-based and super-population targets, and the identifiability
conditions and transportability limits follow Hern\'an and Robins
\cite[ch.~2, p.~13; \S2.1, pp.~14--15; ch.~3, pp.~27--28; \S4.3, p.~52;
Fine Point~4.2, p.~54; \S\S10.1--10.3, pp.~138--144]{HernanRobins2020WhatIf2026};
there, Technical Point~2.1 notes that mean exchangeability suffices for the
equality of the counterfactual mean and the observed conditional mean
\cite[\S2.1, p.~15]{HernanRobins2020WhatIf2026}. The theorems of this chapter,
their proofs and the six-unit table are the handbook's own; they are not
presentations of a particular source.
